% Shared packages, layout, theorem environments, and reusable macros.
% The template is intentionally venue-neutral. Replace the document class and
% citation style in main.tex when a target venue provides its own style file.

% ---- Encoding and typography ------------------------------------------------
% XeLaTeX/LuaLaTeX uses ctex so that the Chinese drafting guidance below is
% rendered correctly. The pdfLaTeX branch keeps the template usable for an
% English-only manuscript.
\usepackage{iftex}
\ifPDFTeX
  \usepackage[T1]{fontenc}
  \usepackage[utf8]{inputenc}
  \usepackage{lmodern}
\else
  % Fandol is shipped with TeX Live and avoids relying on system-specific
  % CJK fonts (for example STHeiti on macOS).
  \usepackage[fontset=fandol]{ctex}
\fi
\usepackage{microtype}
\usepackage{xspace}

% ---- Page layout ------------------------------------------------------------
\usepackage[a4paper,margin=1in]{geometry}
\usepackage{setspace}
\setstretch{1.08}
\usepackage{parskip}
\setlength{\parindent}{1.5em}
\usepackage{placeins}

% ---- Mathematics, algorithms, and units -----------------------------------
\usepackage{amsmath,amssymb,amsthm}
\usepackage{mathtools}
\usepackage{siunitx}
\usepackage{algorithm}
\usepackage{algpseudocode}

% ---- Figures, tables, and lists --------------------------------------------
\usepackage{graphicx}
\graphicspath{{figures/}{figures/main/}{figures/appendix/}}
\usepackage{booktabs}
\usepackage{multirow}
\usepackage{tabularx}
\usepackage{array}
\usepackage{makecell}
\usepackage{subcaption}
\usepackage{float}
\usepackage{enumitem}
\usepackage{xcolor}

% ---- Citations and links ----------------------------------------------------
\usepackage[round,authoryear]{natbib}
\usepackage{xurl}
\usepackage{hyperref}
\usepackage[nameinlink,noabbrev]{cleveref}

\hypersetup{
  colorlinks=true,
  linkcolor=blue!55!black,
  citecolor=green!40!black,
  urlcolor=blue!65!black,
  pdftitle={Your Paper Title}
}

% ---- Colors and writing helpers --------------------------------------------
\definecolor{todoRed}{HTML}{B00020}
\definecolor{noteBlue}{HTML}{0B5394}
\definecolor{lightGray}{HTML}{F3F4F6}

% Visible markers make unfinished sections easy to find during drafting.
% Remove or redefine this command before submission if desired.
\newcommand{\placeholder}[1]{\textcolor{todoRed}{\textbf{[TODO: #1]}}}
\newcommand{\note}[1]{\textcolor{noteBlue}{\textit{[Note: #1]}}}
\newcommand{\methodname}{\textsc{MethodName}\xspace}
\newcommand{\datasetname}{\textsc{DatasetName}\xspace}
\newcommand{\eg}{e.g.,\xspace}
\newcommand{\ie}{i.e.,\xspace}

% A blank figure placeholder that compiles before real assets are available.
\newcommand{\placeholderfigure}[2]{%
  \begin{figure}[t]
    \centering
    \fcolorbox{black!35}{lightGray}{%
      \rule{0pt}{1.35in}\rule{0.88\linewidth}{0pt}%
    }
    \caption{#1}
    \label{#2}
  \end{figure}%
}

% ---- Theorem-like environments --------------------------------------------
\theoremstyle{plain}
\newtheorem{theorem}{Theorem}
\newtheorem{proposition}{Proposition}
\newtheorem{lemma}{Lemma}
\newtheorem{corollary}{Corollary}
\theoremstyle{definition}
\newtheorem{definition}{Definition}
\newtheorem{example}{Example}
\theoremstyle{remark}
\newtheorem{remark}{Remark}

% ---- Table helpers ----------------------------------------------------------
\newcolumntype{Y}{>{\centering\arraybackslash}X}
\newcommand{\best}[1]{\textbf{#1}}
\newcommand{\secondbest}[1]{\underline{#1}}

% ---- Consistent figure/table references ------------------------------------
\crefname{algorithm}{Algorithm}{Algorithms}
\crefname{figure}{Figure}{Figures}
\crefname{table}{Table}{Tables}
\crefname{equation}{Equation}{Equations}
